\section{问题重述}

霍尔木兹海峡连接波斯湾和阿曼湾，是中东原油进入国际市场的重要通道。若该通道发生封锁，全球原油市场将面对突发供应缺口。赛题要求分析封锁对国际原油价格的影响，并特别关注短期价格冲击、缓冲机制、持续 60--180 天时的价格路径以及关键因素敏感性。

本文将问题拆解为四个层次：

\begin{enumerate}
  \item 数据层：清洗附件价格数据，确定真实拟合口径和冲突窗口。
  \item 基准层：用传统低弹性供需模型估计无缓冲理论上界。
  \item 动态层：引入 SPR、库存、绕道、需求调整、恐慌衰减和预期修复，解释现实价格路径。
  \item 外推层：构建 60--180 天三情景预测，并通过敏感性分析识别关键变量。
\end{enumerate}

\begin{table}[H]
\centering
\caption{赛题任务与本文模型回应}
\begin{tabularx}{\linewidth}{lX X}
\toprule
赛题要求 & 题面重点 & 本文回应 \\
\midrule
任务 1 & 建立 0--60 天日度供需平衡模型，使用附件布伦特原油价格数据，并考虑供应中断、SPR、库存、绕道、恐慌需求和短期弹性。 & 完成附件数据清洗、传统供需基准、短期动态递推、参数校准和防御性检验，给出 RMSE、MAE、MAPE、基准对比、滞后检验、机制消融和过拟合压力测试。 \\
任务 2 & 若封锁持续到 60--180 天，预测油价变化和价格平衡点，考虑增产、长期弹性、阿联酋增产约束和库存耗尽跳变风险。 & 构建乐观、中性、悲观三情景路径，给出第 60/90/120/180 天价格、外推期峰值、库存风险、二次跳涨风险和关键参数敏感性排序。 \\
\bottomrule
\end{tabularx}
\end{table}
